%%=============================================================================
%% Methodologie
%%=============================================================================

\chapter{\IfLanguageName{dutch}{Methodologie}{Methodology}}%
\label{ch:methodologie}

%% TODO: In dit hoofstuk geef je een korte toelichting over hoe je te werk bent
%% gegaan. Verdeel je onderzoek in grote fasen, en licht in elke fase toe wat
%% de doelstelling was, welke deliverables daar uit gekomen zijn, en welke
%% onderzoeksmethoden je daarbij toegepast hebt. Verantwoord waarom je
%% op deze manier te werk gegaan bent.
%% 
%% Voorbeelden van zulke fasen zijn: literatuurstudie, opstellen van een
%% requirements-analyse, opstellen long-list (bij vergelijkende studie),
%% selectie van geschikte tools (bij vergelijkende studie, "short-list"),
%% opzetten testopstelling/PoC, uitvoeren testen en verzamelen
%% van resultaten, analyse van resultaten, ...
%%
%% !!!!! LET OP !!!!!
%%
%% Het is uitdrukkelijk NIET de bedoeling dat je het grootste deel van de corpus
%% van je bachelorproef in dit hoofstuk verwerkt! Dit hoofdstuk is eerder een
%% kort overzicht van je plan van aanpak.
%%
%% Maak voor elke fase (behalve het literatuuronderzoek) een NIEUW HOOFDSTUK aan
%% en geef het een gepaste titel.

In de eerste fase van het onderzoek is een analyse uitgevoerd van het huidige onboarding- en off\-boar\-ding\-pro\-ces\-sen binnen het lokale gemeentebestuur van Berlare. Deze analyse is noodzakelijk om inzicht te verkrijgen in welke stappen momenteel worden toegepast binnen de processen en hoeveel tijd deze stappen in beslag nemen. Hiervoor zijn gesprekken gevoerd met betrokken medewerkers om na te gaan welke taken zij uitvoeren en welke systemen daarbij gebruikt worden.

Daarnaast is in deze fase onderzocht welke fouten zich het vaakst voordoen en waar knelpunten optreden binnen het bestaande proces. Op basis van deze bevindingen is een lijst met requirements opgesteld, die geordend is volgens functionaliteit en prioriteit. Deze fase heeft de basis gevormd voor het verdere onderzoek en is afgerond tijdens de tweede week van februari. \newline

In de tweede fase is een literatuurstudie uitgevoerd. Het doel van deze studie is het in kaart brengen van bestaande technologieën, systemen en methodieken die ingezet worden voor de automatisatie van onboardingprocessen. Door het bestuderen van relevante literatuur en bestaande oplossingen is er inzicht verkregen in mogelijke technische benaderingen die toepasbaar zijn binnen de context van een lokaal bestuur. 
Op basis van deze studie is een selectie gemaakt van systemen en technologieën die potentieel een positieve invloed kunnen hebben op het onboardingproces. Deze fase is deels afgerond tegen het einde van februari. Deels omdat dit een fase is die tijdens het verder uitwerken van de proof of concept, steeds opnieuw is gedaan. Bij het vastlopen van het uitwerken, is er uiteraard opnieuw onderzocht. \newline

De derde en meest omvangrijke fase van het onderzoek heeft bestaan uit het uitwerken van een proof of concept. In deze fase zijn de resultaten uit de analyse- en literatuurfase gecombineerd om een zo realistisch en toepasbaar mogelijke oplossing te ontwerpen en te implementeren. Het doel van dit proof of concept was aan te tonen hoe een geautomatiseerd onboarding- en offboardingproces kan functioneren binnen de bestaande IT-infrastructuur van de gemeente Berlare.

Er wordt hierbij gebruik gemaakt van API-koppelingen, die door Power Automate zelf verwerkt worden, waarbij gegevens vanuit Microsoft Forms gekoppeld worden aan verschillende attributen binnen Active Directory. Daarnaast zijn automatisatie- en orkestratietools, Microsoft Power Automate, ingezet om de robuustheid en schaalbaarheid van de oplossing te verhogen. \newline

Na de implementatie van het proof of concept is een evaluatiefase gevolgd. In deze fase is het ontwikkelde systeem geëvalueerd aan de hand van de eerder opgestelde requirements. Daarbij is nagegaan in welke mate het proof of concept voldoet aan de vooropgestelde functionele en technische verwachtingen. Daarnaast is een vergelijking gemaakt tussen de bestaande processen en de geautomatiseerde oplossing, met bijzondere aandacht voor tijdsbesparing, foutreductie en gebruiksgemak. Waar mogelijk is feedback van betrokken medewerkers meegenomen om sterktes en verbeterpunten van de oplossing te identificeren. \newline

Hoewel de verschillende fasen chronologisch zijn voorgesteld, zijn ze iteratief doorlopen. Inzichten die voortkomen uit de proof-of-concept- en evaluatiefase hebben aanleiding gegeven tot een aanvullend literatuuronderzoek of bijkomende bevragingen bij medewerkers. Deze iteratieve aanpak draagt bij aan een eindresultaat dat realistisch, toepasbaar en technisch onderbouwd is.\newline

Figuur \ref{fig:Methodologie} toont een eenvoudige versie van de methodologie.

\begin{figure}[htbp]
    \centering
    \includegraphics[width=1\textwidth]{Methodologie.png}
    \caption[Methodologie]{Visuele voorstelling van de methodologie}
    \label{fig:Methodologie}
\end{figure}
